%% Dokumentenklasse (KOMA-Script) ==========================================
%% scrartcl: Artikel-Klasse, ideal für Seminararbeiten und Zeitschriftenartikel
\documentclass[%
   final,%              % Finales Dokument (vs. draft)
   paper=a4,%           % DIN A4 Papierformat
   pagesize=auto,%      % PDF-Seitengröße automatisch einstellen
   fontsize=10pt,%      % Schriftgröße 10pt (typisch für Zeitschriften)
   twocolumn,%          % Zweispaltiges Layout
   version=last,%       % Neueste KOMA-Script Version
   parskip=half%        % Absatzabstand statt Einzug
]{scrartcl}


%% Spracheinstellung ========================================================
\def\lang{ngerman}

%% Preambel laden ===========================================================
%%% === Überschriften ==========================================================
\KOMAoptions{%
   headings=normal,%          % Normale Überschriften (nicht zu groß)
   numbers=noenddot%          % Kein Punkt nach Abschnittsnummern
}%
\setcounter{secnumdepth}{3}

%%% === Seitenlayout ===========================================================
%%% Das konkrete Seitenlayout wird in der Artikel-Präambel über geometry und
%%% scrlayer-scrpage gesetzt, um Doppelkonfigurationen zu vermeiden.
\KOMAoptions{twoside=false}%

%%% === Table of Contents =====================================================
\setcounter{tocdepth}{2}

%%% ----------------------------------------------------------------------
%%% LaTeX-Präambel für Artikel (zweispaltig, populärwissenschaftlich)
%%% Vereinfachte Version basierend auf der Thesis-Präambel
%%% ----------------------------------------------------------------------

%%% 1. Basiskonfiguration
%%% ----------------------------------------------------------------------
%%% Packages for LaTeX - programming
%
% Define commands that don't eat spaces.
\usepackage{xspace}
% IfThenElse
\usepackage{ifthen}
%%% Doc: ftp://tug.ctan.org/pub/tex-archive/macros/latex/contrib/oberdiek/ifpdf.sty
% command for testing for pdf-creation
\usepackage{ifpdf} %\ifpdf \else \fi

%%% Internal Commands: ----------------------------------------------
\makeatletter
%
\providecommand{\IfPackageLoaded}[2]{\@ifpackageloaded{#1}{#2}{}}
\providecommand{\IfPackageNotLoaded}[2]{\@ifpackageloaded{#1}{}{#2}}
\providecommand{\IfElsePackageLoaded}[3]{\@ifpackageloaded{#1}{#2}{#3}}
%
\newboolean{chapteravailable}%
\setboolean{chapteravailable}{false}%

\ifcsname chapter\endcsname
  \setboolean{chapteravailable}{true}%
\else
  \setboolean{chapteravailable}{false}%
\fi


\providecommand{\IfChapterDefined}[1]{\ifthenelse{\boolean{chapteravailable}}{#1}{}}%
\providecommand{\IfElseChapterDefined}[2]{\ifthenelse{\boolean{chapteravailable}}{#1}{#2}}%

\providecommand{\IfDefined}[2]{%
\ifcsname #1\endcsname
   #2 %
\else
     % do nothing
\fi
}

\providecommand{\IfElseDefined}[3]{%
\ifcsname #1\endcsname
   #2 %
\else
   #3 %
\fi
}

\providecommand{\IfElseUnDefined}[3]{%
\ifcsname #1\endcsname
   #3 %
\else
   #2 %
\fi
}


%
% Check for 'draft' mode - commands.
\newcommand{\IfNotDraft}[1]{\ifx\@draft\@undefined #1 \fi}
\newcommand{\IfNotDraftElse}[2]{\ifx\@draft\@undefined #1 \else #2 \fi}
\newcommand{\IfDraft}[1]{\ifx\@draft\@undefined \else #1 \fi}
%

% Definde frontmatter, mainmatter and backmatter if not defined
\@ifundefined{frontmatter}{%
   \newcommand{\frontmatter}{%
      %In Roemischen Buchstaben nummerieren (i, ii, iii)
      \pagenumbering{roman}
   }
}{}
\@ifundefined{mainmatter}{%
   % scrpage2 benoetigt den folgenden switch
   % wenn \mainmatter definiert ist.
   \newif\if@mainmatter\@mainmattertrue
   \newcommand{\mainmatter}{%
      % -- Seitennummerierung auf Arabische Zahlen zuruecksetzen (1,2,3)
      \pagenumbering{arabic}%
      \setcounter{page}{1}%
   }
}{}
\@ifundefined{backmatter}{%
   \newcommand{\backmatter}{
      %In Roemischen Buchstaben nummerieren (i, ii, iii)
      \pagenumbering{roman}
   }
}{}

% Pakete speichern die spaeter geladen werden sollen
\newcommand{\LoadPackagesNow}{}
\newcommand{\LoadPackageLater}[1]{%
   \g@addto@macro{\LoadPackagesNow}{%
      \usepackage{#1}%
   }%
}

\newcommand*{\printGenerativeAIDeclaration}{%
\ifcsdef{declareUseOfGenerativeAITool}{%
\ifthenelse{\equal{\lang}{ngerman}}%
{%
\section*{Erklärung zur Verwendung von generativer KI}
In Übereinstimmung mit den Empfehlungen der
Deutschen Forschungsgemeinschaft (DFG)\footnote{DFG formuliert Richtlinien für den Umgang mit generativen Modellen für Text- und Bild: \url{https://www.dfg.de/resource/blob/289676/89c03e7a7a8a024093602995974832f9/230921-statement-executive-committee-ki-ai-data.pdf}}
und denen der Zeitschrift Theoretical Computer Science\footnote{Erklärung zur Verwendung von generativer KI in wissenschaftlichen Arbeiten: \url{https://www.sciencedirect.com/journal/theoretical-computer-science/publish/guide-for-authors}}
erkläre ich (der Autor/die Autorin) hiermit den Einsatz von generativer KI.

Bei der Vorbereitung dieser Arbeit habe ich
\declareUseOfGenerativeAITool\ verwendet, um ausschließlich die Lesbarkeit und Sprache zu verbessern.
Nach der Verwendung von \declareUseOfGenerativeAITool\ habe ich
den Inhalt überprüft und nach Bedarf bearbeitet und übernehme die volle Verantwortung
für den Inhalt dieser Arbeit.
}
{%
\section*{Declaration of the use of Generative AI}
In accordance with the recommendation of
the German Research Foundation (DFG - Deutsche Forschungsgemeinschaft)\footnote{DFG Formulates Guidelines for Dealing with Generative Models for Text and Image Creation: \url{https://www.dfg.de/en/news/news-topics/announcements-proposals/2023/info-wissenschaft-23-72}}
and that of the journal Theoretical Computer Science\footnote{Declaration of generative AI in scientific writing: \url{https://www.sciencedirect.com/journal/theoretical-computer-science/publish/guide-for-authors}}
I (the author) herby declare the use of generative AI.

During the preparation of this work I used
\declareUseOfGenerativeAITool\ in order to improve readability and language, only.
After using \declareUseOfGenerativeAITool, I reviewed and edited
the content as needed and take full responsibility
for the content of this thesis.
}
}{}
}

\newcommand*\printGenderDisclaimer{%
\ifthenelse{\equal{\lang}{ngerman}}%
{%
\section*{Gender-Disclaimer}
\thispagestyle{empty}
In dieser Arbeit wird aus Gründen der besseren Lesbarkeit das generische Maskulinum verwendet. Weibliche und anderweitige Geschlechteridentitäten werden dabei ausdrücklich mitgemeint, soweit es für die Aussage erforderlich ist.
}
{%
\section*{Gender Disclaimer}
For reasons of better readability, the generic masculine is used in this work. Female and other gender identities are explicitly included where necessary for the statement.
}
}

\newcommand*\printDeclarationOfIndependence{%
\ifthenelse{\equal{\lang}{ngerman}}%
{%
\section*{Erklärung der Selbstständigkeit}
\thispagestyle{empty}
Hiermit versichere ich, die vorliegende Arbeit selbstständig verfasst und keine anderen als die angegebenen Quellen und Hilfsmittel benutzt sowie die Zitate deutlich kenntlich gemacht zu haben.
\vspace{4\baselineskip}\\
Gießen, den \today \hfill \theAuthor
\vspace{4\baselineskip}\\
}
{%
\section*{Declaration of Independence}
\thispagestyle{empty}
I hereby declare that I have composed the present work independently and have not used any sources or aids other than those cited, and that all quotations have been clearly indicated.
\vspace{4\baselineskip}\\
Gießen, on \today \hfill \theAuthor
\vspace{4\baselineskip}\\
}
}

\let\oldauthor\author
\renewcommand{\author}[1]{\oldauthor{#1}\newcommand{\theAuthor}{#1}}

\let\oldtitle\title
\renewcommand{\title}[1]{\oldtitle{#1}\newcommand{\theTitle}{#1}}

\newcommand*{\academicTitle}[1]{\newcommand{\theAcedemicTitle}{#1}}
\newcommand*{\thesis}[1]{\newcommand{\theThesis}{#1}}
\newcommand*{\firstReferee}[1]{\newcommand{\theFirstReferee}{#1}}
\newcommand*{\secondReferee}[1]{\newcommand{\theSecondReferee}{#1}}
\newcommand*{\studentID}[1]{\newcommand{\theStudentID}{#1}}
\newcommand*{\studentAddress}[1]{\newcommand{\theStudentAddress}{#1}}

\makeatother
%%% ----------------------------------------------------------------

%%% Spracheinstellungen (Babel)
\ifthenelse{\equal{\lang}{ngerman}}
  {\usepackage[ngerman]{babel}}
  {\usepackage[english]{babel}}

%%% Farben
\usepackage[table]{xcolor}

%%% Bilder einbinden
\usepackage{graphicx}

%%% Rechnen mit LaTeX
\usepackage{calc}


%%% 2. Text & Typografie
%%% ----------------------------------------------------------------------

%%% Schrift-Encoding
\usepackage[T1]{fontenc}
\usepackage{textcomp}

%%% Schriftarten laden
%% ========================================================================
%% Schriftarten-Konfiguration
%% ========================================================================
%% Latin Modern ist die Standard-Schriftart für LaTeX-Dokumente.
%% Sie sieht in PDFs ausgezeichnet aus und ist hochwertig gestaltet.
%%
%% Alternative Schriften können durch Auskommentieren aktiviert werden:
%% - Times/Helvetica/Courier: Klassische Kombination
%% - Palantino: Elegant und gut lesbar
%% - Charter: Gute Lesbarkeit bei kleineren Größen
%%
%% WICHTIG: Nicht das 'ae' Paket laden! Das führt zu schlechten Fonts.

%% Aktive Schriftart: Latin Modern (Standard)
\usepackage{lmodern}

%% Alternative 1: Times, Helvetica, Courier (Word-Standard)
%\usepackage{mathptmx}
%\usepackage[scaled=.90]{helvet}
%\usepackage{courier}

%% Alternative 2: Palantino, Helvetica, Courier
%\usepackage{mathpazo}
%\usepackage[scaled=.95]{helvet}
%\usepackage{courier}

%% Alternative 3: Charter mit Bera Sans
%\usepackage{charter}\linespread{1.05}
%\renewcommand{\sfdefault}{fvs}





%%% Zeilenabstand (einzeilig für Zeitschriftenartikel)
\usepackage{setspace}
\singlespacing

%%% Hurenkinder und Schusterjungen vermeiden
\widowpenalty=10000
\clubpenalty=10000

%%% Layout-Anpassung: Kompakte Überschriften für Artikel
\RedeclareSectionCommand[
  beforeskip=-1\baselineskip,
  afterskip=0.25\baselineskip
]{section}
\RedeclareSectionCommand[
  beforeskip=-0.75\baselineskip,
  afterskip=0.15\baselineskip
]{subsection}
\RedeclareSectionCommand[
  beforeskip=-0.5\baselineskip,
  afterskip=0.1\baselineskip
]{subsubsection}

%%% Mikrotypografie
\usepackage[protrusion=true]{microtype}
\emergencystretch=2em

%%% Besserer Flattersatz
\usepackage{ragged2e}

%%% Intelligente Auslassungspunkte
\usepackage{ellipsis}

%%% Textauszeichnungen
\usepackage[normalem]{ulem}
\usepackage{soul}

%%% URLs umbrechen
\usepackage{xurl}

%%% Anführungszeichen
\usepackage[babel, german=quotes]{csquotes}
\SetBlockThreshold{2}


%%% 3. Layout & Seitenränder
%%% ----------------------------------------------------------------------

%%% Seitenränder für zweispaltigen Artikel (etwas großzügiger)
\usepackage[
  a4paper,
  left=20mm,
  right=20mm,
  top=25mm,
  bottom=25mm,
  columnsep=8mm
]{geometry}

%%% Kopf- und Fußzeilen (KOMA-Script) - vereinfacht
\usepackage[automark]{scrlayer-scrpage}
\clearpairofpagestyles
\ohead{\pagemark}
\cfoot{}
\KOMAoptions{headsepline=.4pt}
\setkomafont{pageheadfoot}{\normalfont\small\sffamily}
\setkomafont{pagenumber}{\sffamily}

%%% Fußnoten
\usepackage[bottom, stable, multiple]{footmisc}
\deffootnote{1.5em}{1em}{\makebox[1.5em][l]{\thefootnotemark}}

%%% Hurenkinder und Schusterjungen vermeiden
\clubpenalty = 10000
\widowpenalty = 10000


%%% 4. Tabellen & Abbildungen
%%% ----------------------------------------------------------------------

\usepackage{booktabs}
\usepackage{multirow}
\usepackage{array}
\usepackage{tabularx}
\usepackage{colortbl}
\usepackage{wrapfig}            % Text um Abbildungen fließen lassen

\definecolor{rowcolorgray}{gray}{0.90}
\definecolor{headercolor}{RGB}{120,150,180}
\newcolumntype{Y}{>{\RaggedRight\arraybackslash}X}

\renewcommand{\arraystretch}{1.3}
\setlength{\tabcolsep}{6pt}

%%% Abbildungen & Floats
\usepackage{float}
\usepackage{flafter}
\usepackage[section]{placeins}

%%% Beschriftungen
\usepackage{caption}
\captionsetup{
   margin = 5pt,
   font = {small,rm},
   labelfont = {small,bf},
   format = plain,
   labelsep = colon,
   justification = RaggedRight,
   singlelinecheck = true,
   position = bottom
}
\DeclareCaptionOption{parskip}[]{}
\DeclareCaptionOption{parindent}[]{}

%%% Farbige Boxen
\usepackage[most]{tcolorbox}


%%% 5. Literatur
%%% ----------------------------------------------------------------------

\usepackage[style=numeric, backend=biber, sorting=none]{biblatex}
\addbibresource{bib/BibtexDatabase.bib}


%%% 6. Links
%%% ----------------------------------------------------------------------

\usepackage[
   hidelinks,
   bookmarks=true,
   pdftitle={Wenn KI zum Programmierer wird},
   pdfauthor={Maria Musterfrau},
]{hyperref}


%%% 7. Sonstiges
%%% ----------------------------------------------------------------------

%%% Listen kompakt
\usepackage{enumitem}
\setlist[itemize,enumerate]{nosep, leftmargin=*}
\setlist[description]{nosep, style=unboxed, leftmargin=0pt, labelsep=0.5em, font=\normalfont\bfseries}

%%% Abstract-Umgebung für Artikel
\renewenvironment{abstract}{%
  \begin{tcolorbox}[
    colback=gray!5,
    colframe=gray!50,
    title={\textbf{Zusammenfassung}},
    fonttitle=\sffamily\small,
    boxrule=0.5pt,
    arc=2mm,
    left=3mm, right=3mm, top=2mm, bottom=2mm
  ]
  \small
}{%
  \end{tcolorbox}
  \vspace{1em}
}

%%% Infobox-Umgebung für Kästen
\newtcolorbox{infobox}[1][]{
  colback=blue!5,
  colframe=blue!40,
  title={#1},
  fonttitle=\sffamily\bfseries\small,
  boxrule=0.5pt,
  arc=2mm,
  left=3mm, right=3mm, top=2mm, bottom=2mm
}

%%% Schriftart für Überschriften
\newcommand\SectionFontStyle{\sffamily}
\setkomafont{sectioning}{\SectionFontStyle}
\addtokomafont{sectioning}{\color{black}}

%%% Titel-Anpassung für Artikel
\setkomafont{title}{\LARGE\sffamily\bfseries}
\setkomafont{subtitle}{\large\sffamily}
\setkomafont{author}{\large\sffamily}
\setkomafont{date}{\normalsize\sffamily}

%%% Titel zentriert über beide Spalten
\makeatletter
\renewcommand{\maketitle}{%
  \twocolumn[%
    \begin{@twocolumnfalse}
      \begin{center}
        {\usekomafont{title}\@title\par}
        \vspace{0.5em}
        {\usekomafont{subtitle}\@subtitle\par}
        \vspace{1em}
        {\usekomafont{author}\@author\par}
        \vspace{0.5em}
        {\usekomafont{date}\@date\par}
      \end{center}
      \vspace{1.5em}
    \end{@twocolumnfalse}
  ]%
}
\makeatother


%% Metadaten des Artikels ==================================================
\author{Maria Musterfrau}
\title{Wenn KI selbst programmiert: Wie Agenten die Softwareentwicklung verändern}
\subtitle{Warum die nächste KI-Welle für Entwickler so folgenreich ist}
\date{\today}

%% Dokument Beginn %%%%%%%%%%%%%%%%%%%%%%%%%%%%%%%%%%%%%%%%%%%%%%%%%%%%%%%%

\begin{document}

% Titel
\maketitle

% Abstract (einspaltiger Kasten)
\begin{abstract}
\noindent
Künstliche Intelligenz schreibt längst nicht mehr nur Texte. Sie beginnt, selbstständig an Software mitzuarbeiten: Agenten testen Programme, suchen Fehler, schlagen Korrekturen vor und arbeiten sich schrittweise an Lösungen heran. Dieser Artikel zeigt, wie diese Systeme funktionieren, wo sie schon heute nützlich sind und warum sie den Entwickleralltag grundlegend verändern könnten.
\end{abstract}

% Hauptinhalt
% Populärwissenschaftlicher Artikel über Agentic AI in der Softwareentwicklung
% Zielgruppe: Leserschaft einer Computerzeitschrift (interessierte Menschen mit IT-Affinität)

\section{Die nächste Stufe der KI-Revolution}

ChatGPT formuliert Texte, Copilot ergänzt Code. Doch die nächste Stufe der KI greift tiefer in die Arbeit von Entwicklerinnen und Entwicklern ein: sogenannte \emph{Agenten}~\cite{wang2023survey}. Sie antworten nicht nur, sie handeln. Sie lesen Dateien, starten Tests, suchen im Netz, entwerfen einen Plan und bessern nach, wenn etwas schiefläuft. Aus dem Gesprächspartner wird damit eine digitale Unterstützung auf Probe.

Genau darin liegt die Faszination dieses Trends und sein Konfliktpotenzial zugleich. Was davon ist schon produktive Gegenwart? Und was klingt nur deshalb überzeugend, weil die Erzählung von der autonomen Maschine so verführerisch ist?

\subsection{Vom Chatbot zum Agenten}

Der Unterschied zwischen einem klassischen Chatbot und einem Agenten zeigt sich am besten an einem einfachen Beispiel:

\textbf{Chatbot:} Sie fragen nach einem Python-Skript zum Sortieren einer Liste. Der Bot liefert Code – ob dieser funktioniert, müssen Sie selbst prüfen.

\textbf{Agent:} Sie beschreiben das Problem. Der Agent schreibt den Code, führt ihn aus, stößt auf einen Fehler, sucht nach der Ursache, korrigiert den Code und testet erneut – weitgehend selbstständig.

Der Kern dieser Fähigkeit: Agenten können \emph{nachdenken}, \emph{handeln} und \emph{aus dem Ergebnis lernen}. In der Forschung wird dieses Grundmuster oft mit dem Kürzel \enquote{ReAct} beschrieben~\cite{yao2023react}.


\subsection{Die Evolution der KI-Assistenten}

Um zu verstehen, wohin die Reise geht, lohnt ein Blick zurück. Tabelle~\ref{tab:evolution} zeigt die Entwicklung von einfachen Autovervollständigungen bis hin zu autonomen Agenten.

\begin{table}[ht]
\centering
\footnotesize
\begin{tabularx}{\columnwidth}{@{}Y Y Y@{}}
\toprule
\textbf{Generation} & \textbf{Beispiel} & \textbf{Fähigkeit} \\
\midrule
1 – Autovervollst. & IntelliSense & Wort vorhersagen \\
2 – Generierung & Copilot & Code aus Kommentaren \\
3 – Konversation & ChatGPT & Fragen beantworten \\
4 – Agenten & Devin, Claude & Autonom lösen \\
\bottomrule
\end{tabularx}
\caption{Evolution der KI-Assistenten in der Softwareentwicklung}
\label{tab:evolution}
\end{table}

Der Sprung von Generation~3 zu~4 ist tiefgreifend: Erstmals übernimmt die KI nicht nur einzelne Schritte, sondern ganze Arbeitsabläufe. Sie schreibt nicht mehr bloß mit, sie arbeitet mit.


\section{Wie ein KI-Agent denkt}

Der Kern der Agententechnik lässt sich in vier Schritten zusammenfassen: \textbf{nachdenken} (Welche Information fehlt noch?), \textbf{handeln} (ein Werkzeug ausführen), \textbf{beobachten} (das Ergebnis prüfen) und \textbf{bewerten} (entscheiden, wie es weitergeht). In Fachkreisen wird dieses Muster oft als \enquote{ReAct} bezeichnet.

Dieser Zyklus wiederholt sich, bis die Aufgabe erledigt ist oder der Agent erkennt, dass er Hilfe braucht. Das klingt unspektakulär, ist aber der entscheidende Unterschied zu klassischen Assistenten, die nur auf einzelne Eingaben reagieren und danach wieder verstummen.

\begin{figure}[!t]
\centering
\includegraphics[width=0.7\columnwidth]{images/agent-workflow.png}
\caption{Arbeitszyklus eines KI-Agenten: vom Ziel bis zum Ergebnis}
\label{fig:agent-workflow}
\end{figure}

\subsection{Ein Blick unter die Haube}

Abbildung~\ref{fig:agent-workflow} zeigt den typischen Ablauf eines Agenten bei einer Programmieraufgabe. Der Prozess beginnt mit einer Zielbeschreibung und durchläuft dann mehrere Schleifen aus Planung, Werkzeugnutzung und Überprüfung. Gerade dieses wiederholte Nachjustieren macht den Unterschied zwischen einer hübschen Demo und einem brauchbaren Werkzeug.

Besonders bemerkenswert ist die Fähigkeit zur Selbstkorrektur~\cite{shinn2023reflexion}: Wenn ein Werkzeugaufruf fehlschlägt oder das Ergebnis nicht passt, ändert der Agent seinen Plan. Genau das macht moderne Agenten so interessant. Sie bleiben nicht am ersten Fehler hängen, sondern suchen nach einem zweiten, dritten oder vierten Weg zum Ziel.

\subsection{Das Gedächtnis des Agenten}

Im Gegensatz zu einfachen Chatbots verfügen moderne Agenten über zwei Arten von Gedächtnis:

\textbf{Episodisches Gedächtnis:} Speichert konkrete Ereignisse der aktuellen Sitzung. \enquote{Ich habe gerade den Test gestartet und er ist fehlgeschlagen.} Diese Informationen helfen dem Agenten, Zusammenhänge zu erkennen und aus Fehlern zu lernen.

\textbf{Semantisches Gedächtnis:} Enthält abstraktes Wissen über die Codebasis. \enquote{Die Datenbankverbindung wird in der Datei config.py konfiguriert.} Dieses Wissen wird oft durch sogenannte Vektor-Datenbanken realisiert, die ähnliche Konzepte schnell auffinden können.

\subsection{Chain-of-Thought: Schritt für Schritt denken}

Eine wichtige Technik moderner Agenten ist das \emph{Schritt-für-Schritt-Denken} (\emph{Chain-of-Thought})~\cite{wei2022chain}: Statt vorschnell zu antworten, zerlegt das Modell eine Aufgabe in kleinere Teilschritte.

Der Unterschied kann erstaunlich groß sein. Auf die Frage \enquote{Ist 17 eine Primzahl?} irrt ein Modell ohne diese Strategie womöglich. Mit der Anweisung \enquote{Denke Schritt für Schritt} prüft es systematisch die Teilbarkeit durch 2, 3, 5 und weitere Zahlen und gelangt verlässlicher zur richtigen Antwort.

Für Programmieraufgaben heißt das: Der Agent untersucht zunächst das Problem, sucht die relevanten Dateien, plant sein Vorgehen und greift erst dann ein. Das senkt die Fehlerquote und macht nachvollziehbar, warum das System gerade diesen Schritt wählt. Gerade beim Fehlersuchen ist diese Transparenz Gold wert.


\section{Die Werkzeugkiste des Agenten}

Ein Agent ist nur so gut wie seine Werkzeuge. Moderne Systeme können auf eine ganze Werkzeugkiste zugreifen:

\begin{itemize}
  \item \textbf{Code-Ausführung:} Python, JavaScript oder andere Sprachen direkt ausführen
  \item \textbf{Dateioperationen:} Dateien lesen, schreiben, durchsuchen
  \item \textbf{Webrecherche:} Aktuelle Informationen aus dem Internet abrufen
  \item \textbf{Versionskontrolle:} Git-Befehle ausführen, Änderungen committen
  \item \textbf{Testwerkzeuge:} Automatische Tests starten und Ergebnisse analysieren
  \item \textbf{APIs:} Mit externen Diensten kommunizieren
\end{itemize}

Das Besondere: Der Agent entscheidet selbst, welches Werkzeug wann sinnvoll ist. Er folgt also nicht bloß einem starren Skript, sondern wählt im Kontext aus. Technisch geschieht das oft über strukturierte Ausgaben, die von der Software als Funktionsaufrufe verstanden werden~\cite{schick2023toolformer}.

\subsection{Werkzeug-Schnittstellen im Detail}

Damit ein Agent ein Werkzeug nutzen kann, muss dieses nach einem standardisierten Schema beschrieben werden. Tabelle~\ref{tab:tools} zeigt beispielhaft, wie typische Werkzeuge aus der Softwareentwicklung für Agenten aufbereitet werden.

\begin{table}[ht]
\centering
\footnotesize
\begin{tabularx}{\columnwidth}{@{}Y Y Y@{}}
\toprule
\textbf{Tool} & \textbf{Input} & \textbf{Output} \\
\midrule
run\_tests & Modul & Testergebnis \\
read\_file & Pfad & Dateiinhalt \\
search & Begriff & Codestellen \\
exec\_cmd & Befehl & Ausgabe \\
create\_file & Pfad+Inhalt & OK/Fehler \\
\bottomrule
\end{tabularx}
\caption{Typische Agenten-Werkzeuge}
\label{tab:tools}
\end{table}

Weil diese Schnittstellen standardisiert sind, lassen sich neue Werkzeuge vergleichsweise leicht ergänzen. Der Agent muss nur verstehen, wofür sie gedacht sind und welche Eingaben sie erwarten. Genau so wächst aus einem einfachen Helfer nach und nach ein System, das immer mehr Arbeitsschritte übernehmen kann.

\subsection{Wie entscheidet der Agent?}

Die Werkzeugauswahl ist eine der Kernfähigkeiten eines Agenten. Der Prozess läuft vereinfacht so ab:

\begin{enumerate}
  \item Der Agent analysiert die aktuelle Aufgabe
  \item Er generiert eine Liste möglicher Werkzeuge
  \item Für jedes Werkzeug bewertet er die Erfolgswahrscheinlichkeit
  \item Das vielversprechendste Werkzeug wird aufgerufen
  \item Das Ergebnis fließt in die nächste Entscheidung ein
\end{enumerate}

Dabei kann der Agent auch \emph{Ketten} von Werkzeugen planen: \enquote{Erst suche ich die relevante Datei, dann lese ich sie, dann analysiere ich den Fehler, dann schreibe ich den Fix.}

Diese Fähigkeit zur vorausschauenden Planung unterscheidet leistungsfähige Agenten von einfachen Systemen, die nur reaktiv auf einzelne Anfragen reagieren.


\section{Praxisbeispiel: Fehlersuche im Code}

Stellen Sie sich einen typischen Nachmittag in einem Entwicklungsteam vor: Drei von hundert automatischen Tests schlagen plötzlich fehl. Klassisch würde nun jemand:

\begin{enumerate}
  \item Die Fehlermeldungen analysieren
  \item Den betroffenen Code finden
  \item Die Ursache verstehen
  \item Eine Lösung entwickeln
  \item Testen, ob die Lösung funktioniert
\end{enumerate}

Ein KI-Agent durchläuft im Grunde dieselben Schritte, nur automatisch und in hoher Geschwindigkeit:

\begin{quote}
\textit{\enquote{Ich sehe 3 fehlgeschlagene Tests in der Datei test\_login.py. Lass mich die Fehlermeldung genauer anschauen... Es scheint ein Problem mit der Datenbankverbindung zu sein. Ich überprüfe die Konfiguration... Hier fehlt ein Timeout-Parameter. Ich füge ihn hinzu und starte die Tests erneut... Alle Tests bestanden!}}
\end{quote}

\subsection{Erfolgsraten in der Praxis}

Wie gut funktionieren diese Systeme wirklich? Die Forschung prüft das mit standardisierten Vergleichstests. Historisch besonders wichtig war \emph{SWE-bench}~\cite{jimenez2023swe}, ein Datensatz mit über 2.000 echten Programmierproblemen aus GitHub-Projekten. Systeme wie SWE-agent~\cite{yang2024sweagent} und Agentless~\cite{zhang2024agentless} haben gezeigt, dass eine kluge Verbindung aus Modell, Werkzeugen und Ausführungsumgebung die Erfolgsraten deutlich steigern kann.

Im Frühjahr 2026 ist die Lage allerdings unübersichtlicher als noch vor einem Jahr: Ranglisten, Modellnamen und Testverfahren ändern sich im Monatsrhythmus. Außerdem gilt der lange populäre Teilbenchmark \enquote{SWE-bench Verified} nur noch eingeschränkt als verlässlicher Maßstab, weil Testdesign und mögliche Trainingskontamination die Aussagekraft verzerren können. Deshalb rücken robustere Setups wie \enquote{SWE-bench Pro} und stärker ganzheitliche Bewertungen in den Vordergrund.~\cite{openai2026swebenchverified,swebench2026leaderboard}

\begin{table}[ht]
\centering
\small
\begin{tabularx}{\columnwidth}{@{}Y c c@{}}
\toprule
\textbf{Phase / Beispiel} & \textbf{Signal} & \textbf{Jahr} \\
\midrule
Einfache Prompts ohne Werkzeuge & 3--5\,\% & 2023 \\
SWE-agent & 12,5\,\% & 2024 \\
Agentless & 27\,\% & 2024 \\
Leistungsstarke Agenten mit Werkzeugen & rund 50\,\% & 2025 \\
Neue Pro-Benchmarks & kein fixer Einzelwert & 2026 \\
\bottomrule
\end{tabularx}
\caption{Historische Meilensteine der Agenten-Evaluation bis Frühjahr 2026}
\label{tab:benchmark}
\end{table}

Die Zahlen in Tabelle~\ref{tab:benchmark} zeigen, wie schnell sich das Feld bewegt: Innerhalb von zwei Jahren wurden aus experimentellen Vorführungen ernstzunehmende Werkzeuge für Entwicklerinnen und Entwickler. Gleichzeitig gilt 2026 mehr denn je: Einzelne Prozentwerte sind nur dann aussagekräftig, wenn man auch Testaufbau, Werkzeuge und Ausführungsumgebung kennt. Gerade bei komplexen Problemen mit vielen Abhängigkeiten stoßen aktuelle Systeme weiterhin an Grenzen. Der Fortschritt ist real, aber er kommt nicht als geradlinige Erfolgsgeschichte daher.

\subsection{Welche Aufgaben eignen sich?}

Nicht jede Programmieraufgabe ist gleich gut für Agenten geeignet. Tabelle~\ref{tab:eignung} gibt eine Orientierung:

\begin{table}[ht]
\centering
\small
\begin{tabularx}{\columnwidth}{@{}Y c c@{}}
\toprule
\textbf{Aufgabentyp} & \textbf{Eignung} & \textbf{Erfolgsrate} \\
\midrule
Lint-Fehler beheben & Sehr gut & 80--90\,\% \\
Unit-Tests schreiben & Gut & 70--80\,\% \\
Einfache Fehlerkorrekturen & Gut & 60--75\,\% \\
Code umstrukturieren & Mittel & 50--70\,\% \\
Neue Funktionen umsetzen & Mittel & 40--60\,\% \\
Software-Architektur entwerfen & Gering & 20--30\,\% \\
\bottomrule
\end{tabularx}
\caption{Eignung verschiedener Aufgaben für KI-Agenten}
\label{tab:eignung}
\end{table}

Faustregel: Je klarer definiert und begrenzter eine Aufgabe, desto besser die Erfolgsaussichten. Offene, kreative oder politisch heikle Probleme bleiben vorerst in menschlicher Hand.

\subsection{Ein konkretes Beispiel aus der Praxis}

Um die Arbeitsweise eines Agenten greifbar zu machen, hilft ein protokollierter Durchlauf bei einer typischen Fehlerkorrektur:

\begin{quote}
\footnotesize
\textbf{Aufgabe:} \enquote{Fix failing test in user\_auth}

\textbf{1. Denken:} \enquote{Welcher Test schlägt fehl?}\\
\textbf{2. Aktion:} \texttt{run\_tests("user\_auth")}\\
\textbf{3. Beobachtung:} \enquote{test\_login\_timeout: Timeout}\\
\textbf{4. Denken:} \enquote{Timeout-Problem. Konfiguration?}\\
\textbf{5. Aktion:} \texttt{read\_file("config/db.py")}\\
\textbf{6. Beobachtung:} \texttt{timeout = 5}\\
\textbf{7. Denken:} \enquote{Zu niedrig für Tests.}\\
\textbf{8. Aktion:} \texttt{edit\_file(..., timeout=30)}\\
\textbf{9. Aktion:} \texttt{run\_tests("user\_auth")}\\
\textbf{10. Beobachtung:} \enquote{15 Tests bestanden}\\
\textbf{Fazit:} Timeout war zu niedrig konfiguriert.
\end{quote}

Dieser protokollierte Ablauf zeigt die charakteristische Abfolge aus Denken, Handeln und Beobachten. Bemerkenswert ist die Fähigkeit, aus Fehlermeldungen auf mögliche Ursachen zu schließen und gezielt nach Lösungen zu suchen.


\section{Architektur moderner Agenten}

Hinter jedem leistungsfähigen Agenten steckt eine durchdachte Architektur. Die wichtigsten Komponenten sind:

\begin{description}
  \item[Agentensteuerung:] Das \enquote{Gehirn} des Systems. Hier laufen Planung und Entscheidungsfindung zusammen. Diese Komponente ruft das Sprachmodell auf und deutet dessen Antworten.
  
  \item[Werkzeugverzeichnis:] Eine Übersicht aller verfügbaren Werkzeuge mit ihren Beschreibungen und Parametern. Daraus wählt der Agent für jede Situation das passende aus.
  
  \item[Gedächtnisverwaltung:] Verwaltet Kurz- und Langzeitgedächtnis. Sie entscheidet, welche Informationen gespeichert und welche wieder verworfen werden.
  
  \item[Sicherheitsschicht:] Prüft jeden Werkzeugaufruf auf Zulässigkeit. Sie verhindert gefährliche Operationen und protokolliert alle Aktionen.
\end{description}

Diese modulare Struktur~\cite{wu2023autogen} macht es möglich, einzelne Bausteine auszutauschen oder zu verbessern, ohne das ganze System neu zu bauen. In der Praxis heißt das: Ein Unternehmen kann mit einem einfachen Agenten beginnen und ihn Schritt für Schritt erweitern. So entstehen auch spezialisierte Helfer, etwa für Benutzeroberflächen, Datenbanken oder Sicherheitsprüfungen.

\subsection{Vergleich: Cloud vs. Lokal}

Agenten können entweder in der Cloud oder direkt auf dem eigenen Rechner laufen. Beide Ansätze haben Vor- und Nachteile:

\begin{table}[ht]
\centering
\footnotesize
\begin{tabularx}{\columnwidth}{@{}Y c c@{}}
\toprule
\textbf{Aspekt} & \textbf{Cloud} & \textbf{Lokal} \\
\midrule
Leistung & Hoch & Begrenzt \\
Kosten & Nutzung nach Verbrauch & Hardware \\
Datenschutz & Transfer & Lokal \\
Latenz & Netzwerk & Schnell \\
Offline & Nein & Ja \\
Wartung & Anbieter & Selbst \\
\bottomrule
\end{tabularx}
\caption{Cloud vs. lokale Agenten}
\label{tab:cloudlokal}
\end{table}

Für sensible Projekte oder Unternehmen mit strengen Vorgaben beim Datenschutz werden lokale Lösungen immer attraktiver. Modelle wie Llama 4~\cite{meta2025llama4}, Mistral- oder DeepSeek-Varianten ermöglichen schon heute leistungsfähige Agenten auf eigener Hardware, wenn auch oft mit hohem Rechenbedarf. Die Entwicklung schreitet jedoch schnell voran: Was heute noch einen starken Server verlangt, könnte in wenigen Jahren auf einem gewöhnlichen Laptop laufen. Zunehmend beliebt sind auch Mischformen, bei denen heikle Schritte lokal und besonders rechenintensive Aufgaben in der Cloud ausgeführt werden.

\subsection{Sicherheit: Der Agent in der Sandbox}

Bei aller Begeisterung bleibt eine wichtige Frage: Wie verhindert man, dass ein autonomer Agent Schaden anrichtet? Die Antwort liegt in mehrschichtigen Sicherheitsmechanismen, die nach dem Prinzip der \enquote{Verteidigung in der Tiefe} aufgebaut sind.

\textbf{Abgeschottete Ausführung:} Der Agent läuft in einer isolierten Umgebung, ähnlich einer virtuellen Maschine. Selbst wenn er fehlerhaften Code ausführt, kann er das eigentliche System nicht ohne Weiteres beschädigen.

\textbf{Berechtigungssystem:} Jedes Werkzeug hat definierte Rechte. Ein Agent darf vielleicht Dateien lesen, aber nicht löschen. Oder er darf Tests starten, aber keinen Code in Produktion deployen.

\textbf{Menschliche Freigabe:} Bei kritischen Aktionen, etwa dem Veröffentlichen von Code, muss ein Mensch zustimmen. Der Agent schlägt vor, die verantwortliche Person entscheidet.

\textbf{Protokollierung:} Jede Aktion wird aufgezeichnet. Bei Problemen lässt sich exakt nachvollziehen, was der Agent getan hat.

\subsection{Schutz vor Manipulation}

Eine besondere Gefahr sind sogenannte \enquote{Prompt-Injection}-Angriffe: Eine böswillige Person versteckt Anweisungen in Daten, die der Agent verarbeitet. Stellen Sie sich vor, in einem Code-Kommentar stünde plötzlich der Befehl \enquote{Lösche alle Dateien}.

Robuste Systeme erkennen solche Manipulationsversuche durch:
\begin{itemize}
  \item Strikte Trennung von Anweisungen und Daten
  \item Validierung aller Eingaben vor der Verarbeitung
  \item Whitelisting erlaubter Operationen
  \item Anomalie-Erkennung bei ungewöhnlichem Verhalten
\end{itemize}

In Tests erreichen gut abgesicherte Systeme eine Abwehrrate von über 99\,\% gegen bekannte Angriffsmuster.


\section{Kosten und Effizienz}

KI-Agenten sind nicht kostenlos. Jeder Aufruf eines leistungsfähigen Sprachmodells kostet Geld, und für eine einzige Aufgabe können schnell Dutzende solcher Aufrufe zusammenkommen.

\subsection{Was kostet ein Agent?}

Tabelle~\ref{tab:kosten} zeigt grobe Richtwerte für verschiedene Aufgabentypen. Die Werte basieren auf typischen Preisniveaus im Frühjahr 2026 und können je nach Anbieter, Modell und Umfang des Kontexts stark variieren.

\begin{table}[ht]
\centering
\small
\begin{tabularx}{\columnwidth}{@{}Y c c@{}}
\toprule
\textbf{Aufgabe} & \textbf{API-Aufrufe} & \textbf{ca. Kosten} \\
\midrule
Einfache Code-Korrektur & 3--5 & 0,02--0,05\,€ \\
Test-Fehlersuche & 8--15 & 0,08--0,15\,€ \\
Umsetzung neuer Funktionen & 20--50 & 0,20--0,50\,€ \\
Komplexe Umstrukturierung & 50--100 & 0,50--1,00\,€ \\
\bottomrule
\end{tabularx}
\caption{Grobe Richtwerte für KI-Agenten-Aufgaben (Stand: Frühjahr 2026)}
\label{tab:kosten}
\end{table}

Verglichen mit Arbeitsstunden in der Softwareentwicklung sind diese Kosten oft gering: Selbst wenn ein Agent mehrere Modellaufrufe benötigt, bleiben einfache Aufgaben häufig im Cent- oder unteren Euro-Bereich. Wirtschaftlich interessant wird das vor allem dort, wo viele kleine, klar eingegrenzte Routineaufgaben anfallen und sich Zeitgewinne Tag für Tag summieren.

\subsection{Token-Optimierung}

Der größte Kostenfaktor ist der \enquote{Kontext} – also alle Informationen, die der Agent für seine Arbeit benötigt. Auch wenn moderne Modelle sehr lange Eingaben verarbeiten können, sollte eine große Codebasis mit Hunderttausenden Zeilen nicht blind komplett an das Sprachmodell übergeben werden.

Clevere Agenten nutzen daher Techniken wie:
\begin{itemize}
  \item \textbf{Chunking:} Code wird in kleinere Abschnitte unterteilt
  \item \textbf{Retrieval:} Nur relevante Dateien werden geladen
  \item \textbf{Zusammenfassung:} Lange Ausgaben werden komprimiert
\end{itemize}

Diese Optimierungen können die Kosten um 40--60\,\% reduzieren, ohne die Qualität wesentlich zu beeinträchtigen.


\section{Grenzen der Technologie}

So beeindruckend die Fortschritte sind – KI-Agenten sind keine Wunderlösung:

\textbf{Große Projekte:} Bei Millionen Zeilen Code, etwa im Linux-Kernel, stoßen aktuelle Agenten an ihre Grenzen. Der Kontext ist schlicht zu groß, um ihn vollständig zu erfassen.

\textbf{Kreative Aufgaben:} Neue Algorithmen erfinden oder revolutionäre Architekturen entwerfen – das bleibt vorerst menschliche Domäne. Agenten sind gut im Anwenden bekannter Muster, weniger im Erfinden neuer.

\textbf{Zwischenmenschliches:} Anforderungen verstehen, mit Auftraggebern und Teams kommunizieren, Prioritäten setzen: Hier ist menschliche Intelligenz unersetzlich. Ein Agent kann technisch perfekten Code liefern, der am eigentlichen Bedarf vorbeigeht.

\textbf{Halluzinationen:} Sprachmodelle können \enquote{halluzinieren}, also plausibel klingende, aber falsche Informationen erzeugen. Bei Agenten kann das zu schwer auffindbaren Fehlern führen.

\begin{infobox}[Wann Agenten scheitern]
Typische Fälle, in denen aktuelle Agenten an ihre Grenzen stoßen:
\begin{itemize}
  \item Altsysteme ohne Dokumentation
  \item Domänenspezifisches Fachwissen (Medizin, Recht)
  \item Aufgaben mit mehrdeutigen Anforderungen
  \item Systeme mit komplexen Abhängigkeiten
\end{itemize}
\end{infobox}


\section{Die Zukunft: Menschen und Maschinen}

Die spannendste Frage ist deshalb nicht \enquote{Ersetzt KI die Menschen, die programmieren?}, sondern \enquote{Wie arbeiten beide zusammen?}

\subsection{Hybride Zusammenarbeit}

Die Praxis zeigt: Am stärksten sind Menschen und Maschinen im Team. Der Agent übernimmt repetitive Aufgaben, etwa Fehler korrigieren, Code formatieren oder Tests schreiben. Menschen konzentrieren sich auf das Wesentliche: Architekturentscheidungen, Nutzererlebnis und kreative Problemlösung.

Ein typischer Ablauf könnte so aussehen:
\begin{enumerate}
  \item Eine Person aus dem Team beschreibt die neue Funktion
  \item Agent erstellt einen ersten Entwurf samt Tests
  \item Eine Person aus dem Team prüft den Entwurf und gibt Rückmeldung
  \item Agent überarbeitet die Lösung anhand dieser Rückmeldung
  \item Eine Person aus dem Team nimmt finale Anpassungen vor
  \item Agent führt abschließende Qualitätsprüfungen durch
\end{enumerate}

\subsection{Neue Berufsbilder}

Rund um diese Werkzeuge entstehen bereits neue Rollen:

\textbf{Agenten-Training:} Spezialisierung von Agenten auf bestimmte Fachgebiete durch Feintuning und geschickte Aufgabenformulierung.

\textbf{Agenten-Koordination:} Entwurf von Abläufen, in denen mehrere Agenten zusammenarbeiten.

\textbf{KI-Prüfung:} Kontrolle der Sicherheit und Zuverlässigkeit von Agentensystemen.

Erste Unternehmen setzen solche Systeme bereits produktiv ein. GitHub Copilot, Amazon Q Developer und Gemini Code Assist zeigen, wie stark große Technologiekonzerne inzwischen auf solche Entwicklungswerkzeuge setzen.~\cite{github2026copilot,aws2026qdeveloper,google2026geminicodeassist}


\section{Praktische Tipps: So starten Sie}

Möchten Sie selbst mit KI-Agenten experimentieren? Hier einige Empfehlungen für den Einstieg:

\subsection{Einfache Werkzeuge zum Ausprobieren}

Für erste Experimente brauchen Sie keine komplexe Infrastruktur:

\textbf{GitHub Copilot:} Der bekannteste KI-Assistent für Entwicklerinnen und Entwickler. Integriert sich nahtlos in VS Code und andere IDEs. Chat-, Edit- und Agentenfunktionen machen ihn 2026 deutlich mehr als nur ein Autovervollständigungswerkzeug.~\cite{github2026copilot}

\textbf{Cursor:} Ein spezialisierter Editor auf Basis von VS Code mit enger KI-Integration. Besonders geeignet für Experimente mit kontextbezogener Code-Erzeugung und automatischen Umstrukturierungen.

\textbf{ChatGPT oder Claude mit Werkzeugunterstützung:} Die Oberflächen großer Sprachmodelle bieten inzwischen brauchbare Agentenfunktionen. Code lässt sich ausführen, Dateien können hochgeladen und analysiert werden, und mehrstufige Aufgaben werden teilweise selbstständig abgearbeitet.

\subsection{Für Fortgeschrittene: Frei verfügbare Werkzeuge}

Wer tiefer einsteigen möchte, findet in der frei zugänglichen Community rund um Softwareentwicklung leistungsfähige Werkzeuge:

\begin{table}[ht]
\centering
\small
\begin{tabularx}{\columnwidth}{@{}l Y@{}}
\toprule
\textbf{Werkzeug} & \textbf{Stärken} \\
\midrule
LangChain & Größtes Ökosystem, viele Integrationen \\
AutoGPT & Autonome Aufgabenbearbeitung \\
CrewAI & Zusammenarbeit mehrerer Agenten \\
LangGraph & Zustandsbasierte Abläufe \\
\bottomrule
\end{tabularx}
\caption{Populäre frei verfügbare Werkzeuge für Agentenentwicklung}
\label{tab:frameworks}
\end{table}

\subsection{Best Practices}

Aus der Praxis haben sich einige Empfehlungen herauskristallisiert:

\begin{enumerate}
  \item \textbf{Klein anfangen:} Starten Sie mit klar definierten, begrenzten Aufgaben. Komplexe Projekte überfordern aktuelle Agenten schnell.
  
  \item \textbf{Immer prüfen:} Vertrauen Sie KI-generiertem Code nie blind. Sorgfältige Prüfung bleibt unverzichtbar, erst recht bei automatisch erzeugten Änderungen.
  
  \item \textbf{Tests sind Pflicht:} Automatische Tests sind die beste Absicherung gegen fehlerhafte Agenten-Outputs. Je höher die Testabdeckung, desto sicherer der Einsatz.
  
  \item \textbf{Kosten überwachen:} API-Aufrufe können sich schnell summieren. Setzen Sie Budgetgrenzen und überwachen Sie den Verbrauch.
  
  \item \textbf{Dokumentieren:} Halten Sie fest, welche Teile Ihres Codes von Agenten generiert wurden. Das erleichtert spätere Wartung.
\end{enumerate}


\section{Ethische Fragen und Verantwortung}

Mit der Verbreitung von KI-Agenten entstehen auch neue ethische Herausforderungen, die Entwicklerinnen und Entwickler ebenso wie Unternehmen adressieren müssen.

\subsection{Urheberrecht und Lizenzfragen}

Sprachmodelle wurden auf Milliarden von Codezeilen trainiert, darunter auch urheberrechtlich geschützter Code. Wenn ein Agent Code erzeugt, der einem existierenden Projekt stark ähnelt, stellt sich zwangsläufig die Frage: Ist das schon eine Urheberrechtsverletzung?

Erste Gerichtsverfahren laufen bereits. Bis Rechtssicherheit besteht, empfiehlt sich:
\begin{itemize}
  \item Generierten Code auf Ähnlichkeiten prüfen
  \item Lizenzkompatibilität sicherstellen
  \item Im Zweifel manuell umschreiben
\end{itemize}

\subsection{Verantwortung für Fehler}

Wer haftet, wenn KI-generierter Code einen Fehler enthält, der Schaden verursacht? Die Person, die den Agenten eingesetzt hat? Das Unternehmen, das den Agenten entwickelt hat? Oder die Firma, deren Produkt betroffen ist?

Aktuell liegt die Verantwortung bei denjenigen, die KI-generierten Code in Produktion bringen: Wer solche Systeme einsetzt, muss ihre Ergebnisse auch verantworten. Das unterstreicht die Notwendigkeit gründlicher Prüfungen und Tests.

\subsection{Arbeitsplätze und Qualifikation}

Die Sorge, KI könnte Entwicklerinnen und Entwickler ersetzen, ist weit verbreitet. Die Realität ist differenzierter:

\textbf{Routineaufgaben} werden zunehmend automatisiert. Einfache Fehlerkorrekturen, Standardbausteine und wiederkehrende Tests übernehmen Agenten schon heute.

\textbf{Komplexe Aufgaben} bleiben menschlich. Architekturentscheidungen, Nutzerforschung, Teamführung, kreative Problemlösung – diese Fähigkeiten werden eher wichtiger.

\textbf{Neue Rollen} entstehen. Wer Agenten effektiv einsetzen, trainieren und überwachen kann, wird gefragt sein.

Die sinnvollste Haltung lautet daher: Verstehen Sie KI-Werkzeuge als Verstärker Ihrer eigenen Fähigkeiten, nicht als Ersatz.


\section{Ein Blick in die Kristallkugel}

Wohin entwickelt sich die Technologie in den nächsten Jahren? Einige Trends zeichnen sich ab:

\subsection{Multi-Agenten-Systeme}

Statt eines einzelnen Agenten arbeiten mehrere spezialisierte Agenten zusammen: Einer entwirft die Struktur, ein anderer schreibt Code, ein dritter prüft die Qualität, ein weiterer gibt Rückmeldung. Das Projekt \emph{MetaGPT}~\cite{hong2023metagpt} verfolgt genau diesen Ansatz und simuliert ein komplettes Software-Team mit klar verteilten Rollen. Diese Arbeitsteilung verspricht bessere Ergebnisse bei komplexen Aufgaben.

\subsection{Längerer Kontext}

Spitzenmodelle verarbeiten 2026 deutlich mehr als 100.000 Tokens; einige offene Modelle werben sogar mit Kontextfenstern im Millionenbereich.~\cite{meta2025llama4} Dadurch können Agenten große Dokumentationen, viele Dateien oder ganze Codebereiche auf einmal überblicken. In der Praxis bleibt ein kluger Umgang mit dem Kontext trotzdem entscheidend, weil lange Eingaben teuer sind und nicht jedes Modell mit extremen Längen gleich zuverlässig umgeht.

\subsection{Spezialisierte Modelle}

Statt Allzweck-Modelle werden für Softwareentwicklung optimierte Modelle dominant. Diese verstehen Programmierkonzepte tiefer und machen weniger Fehler bei technischen Aufgaben.

\subsection{Integration in Entwicklungsumgebungen}

Die Grenzen zwischen Entwicklungsumgebung und Agent verschwimmen. Künftige Werkzeuge werden Agentenfunktionen nahtlos einbauen: automatische Umstrukturierungen, kontinuierliche Code-Analyse und proaktive Verbesserungsvorschläge.


\section{Fazit: Revolution in Zeitlupe}

KI-Agenten in der Softwareentwicklung sind keine Science-Fiction mehr. Sie sind da. Nicht perfekt, nicht allwissend, aber schon heute nützlich.

Gerade darin liegt ihre eigentliche Sprengkraft. Nicht weil sie Entwicklerinnen und Entwickler morgen überflüssig machen würden, sondern weil sie den Takt der Arbeit verändern: Was gestern mühsame Kleinarbeit war, wird morgen vielleicht zur automatisierten Zwischenstufe.

Die Technologie entwickelt sich rasant. Was heute bei klar abgegrenzten, gut testbaren Aufgaben bereits erstaunlich gut funktioniert, scheitert in offenen, schlecht dokumentierten Projekten oft noch überraschend schnell. Gleichzeitig entstehen neue Fragen: Wer trägt die Verantwortung für Fehler im KI-generierten Code? Wie schützen wir uns vor bösartigen Agenten? Und wie verändert sich die Arbeit von Entwicklerinnen und Entwicklern?

Für Entwicklerinnen und Entwickler bedeutet das: Wer diese Werkzeuge versteht und klug einsetzt, gewinnt Tempo, Reichweite und Spielraum. Wer sie ignoriert, läuft Gefahr, den Anschluss zu verlieren.

Die eigentliche Revolution besteht also nicht darin, dass Maschinen plötzlich allein programmieren. Sie besteht darin, dass Softwareentwicklung zu einer Zusammenarbeit wird, in der Menschen die Richtung vorgeben und Maschinen immer mehr Zwischenschritte übernehmen. Wer das früh versteht, ist für die nächste Phase des digitalen Wandels besser gerüstet.

\begin{infobox}[Zum Weiterlesen]
\begin{itemize}
  \item \textbf{ReAct-Paper:} Das wissenschaftliche Fundament der Agenten-Technologie~\cite{yao2023react}
  \item \textbf{SWE-bench:} Der historische Referenztest für Code-Agenten; für aktuelle Vergleiche lohnt der Blick auf neuere Ranglisten und Pro-Setups~\cite{jimenez2023swe,swebench2026leaderboard,openai2026swebenchverified}
  \item \textbf{LangChain:} Weit verbreitetes frei verfügbares Werkzeug für Agentenentwicklung~\cite{chase2023langchain}
\end{itemize}
\end{infobox}


\section*{Glossar}

\begin{description}
  \item[Agent:] Ein KI-System, das eigenständig Ziele verfolgt, Pläne erstellt und Werkzeuge nutzt – im Gegensatz zu passiven Chatbots.
  
  \item[Chain-of-Thought:] Technik, bei der ein Sprachmodell ein Problem Schritt für Schritt durchgeht, was oft zu besseren Ergebnissen führt.
  
  \item[Halluzination:] Das Phänomen, dass Sprachmodelle plausibel klingende, aber faktisch falsche Informationen generieren.
  
  \item[LLM (Large Language Model):] Großes Sprachmodell wie GPT-4.1, Claude oder Gemini, das auf riesigen Textmengen trainiert wurde.
  
  \item[Prompt Injection:] Angriffstechnik, bei der schädliche Anweisungen in Eingabedaten versteckt werden, um ein KI-System zu manipulieren.
  
  \item[RAG (Retrieval-Augmented Generation):] Methode, bei der ein Sprachmodell vor der Antwortgenerierung relevante Informationen aus einer Datenbank abruft.
  
  \item[ReAct:] Methodik, bei der Agenten bewusst zwischen Nachdenken und Handeln wechseln.
  
  \item[Sandbox:] Isolierte Ausführungsumgebung, die verhindert, dass Code das umgebende System beeinflusst.
  
  \item[Token:] Grundeinheit der Textverarbeitung in Sprachmodellen – grob entspricht ein Token 0,75 Wörtern.
  
  \item[Vektor-Datenbank:] Spezielle Datenbank zur Speicherung und schnellen Suche von semantisch ähnlichen Texten oder Code.
\end{description}


%% Literaturverzeichnis
\printbibliography[title=Quellen und weiterführende Literatur]

% KI-Erklärung als Fußnote am Ende
\vspace{1em}
\noindent\rule{\columnwidth}{0.4pt}
{\footnotesize\textbf{Hinweis zur KI-Nutzung:} Bei der Erstellung dieses Artikels wurden KI-Werkzeuge wie GitHub Copilot und ChatGPT zur Unterstützung bei Recherche und Formulierung eingesetzt. Die inhaltliche Verantwortung liegt bei der Autorin.}

\end{document}
